%!TEX root = ../sztuthesis_main.tex
% 附录示例：正式撰写时可删除本文件或清空后替换为真实附录；多个附录可仿照增加「附录2」等。

\clearpage
\markright{附录1（示例）}

\phantomsection
\addcontentsline{toc}{section}{附录1（示例）}

{\centering \zihao{-2} \bfseries \heiti 附录1（示例）\par}

\vspace{0.5cm}

\zihao{-4} \songti
\setlength{\parindent}{2em}

本段为模版性附录正文示例。在实际论文中，可将不宜在正文展开的长篇公式推导、重复性数据表、主要程序清单等置于附录；图表与公式的编号规则若与正文分章编号不同，请按学院细则在导言区或本文件内另行设定。

本模板由 \cls{SZTUthesis.cls} 提供版式。

参考文献数据写入 \bib{thesis-references.bib}。正文中的图、表和代码示例通常分别放在 \env{figure}、\env{table} 和 \env{verbatim} 环境中。

若使用 \app{VS Code} 与 \app{LaTeX Workshop}，可执行 \oper{Build LaTeX project} 操作完成编译；引用格式相关功能由 \pkg{gbt7714} 支持。
